\documentclass[10pt,a4paper]{amsart}
\usepackage{amsmath,amssymb,amsthm,graphicx}
% Retained minimal preamble for the Stacks Project chapter varieties.tex.
% Verbatim excerpt of the chapter preamble: Xy-pic setup (native preamble lines 14-18), the native theorem
% declarations with the mathematical macros (native preamble lines 154-217), and the project own no-op
% definitions of the reference/slogan/history display environments (native wrapper.tex lines 1-3) that the
% counted statement uses for its slogan. Copied mechanically; nothing added, removed or redefined.
% For commutative diagrams we use Xy-pic
\usepackage[all]{xy}

% We use 2cell for 2-commutative diagrams.
\xyoption{2cell}
% Theorem environments.
%
\theoremstyle{plain}
\newtheorem{theorem}[subsection]{Theorem}
\newtheorem{proposition}[subsection]{Proposition}
\newtheorem{lemma}[subsection]{Lemma}

\theoremstyle{definition}
\newtheorem{definition}[subsection]{Definition}
\newtheorem{example}[subsection]{Example}
\newtheorem{exercise}[subsection]{Exercise}
\newtheorem{situation}[subsection]{Situation}

\theoremstyle{remark}
\newtheorem{remark}[subsection]{Remark}
\newtheorem{remarks}[subsection]{Remarks}

\numberwithin{equation}{subsection}

% Macros
%
\def\lim{\mathop{\mathrm{lim}}\nolimits}
\def\colim{\mathop{\mathrm{colim}}\nolimits}
\def\Spec{\mathop{\mathrm{Spec}}}
\def\Hom{\mathop{\mathrm{Hom}}\nolimits}
\def\Ext{\mathop{\mathrm{Ext}}\nolimits}
\def\SheafHom{\mathop{\mathcal{H}\!\mathit{om}}\nolimits}
\def\SheafExt{\mathop{\mathcal{E}\!\mathit{xt}}\nolimits}
\def\Sch{\mathit{Sch}}
\def\Mor{\mathop{\mathrm{Mor}}\nolimits}
\def\Ob{\mathop{\mathrm{Ob}}\nolimits}
\def\Sh{\mathop{\mathit{Sh}}\nolimits}
\def\NL{\mathop{N\!L}\nolimits}
\def\CH{\mathop{\mathrm{CH}}\nolimits}
\def\proetale{{pro\text{-}\acute{e}tale}}
\def\etale{{\acute{e}tale}}
\def\QCoh{\mathit{QCoh}}
\def\Ker{\mathop{\mathrm{Ker}}}
\def\Im{\mathop{\mathrm{Im}}}
\def\Coker{\mathop{\mathrm{Coker}}}
\def\Coim{\mathop{\mathrm{Coim}}}

% Boxtimes
%
\DeclareMathSymbol{\boxtimes}{\mathbin}{AMSa}{"02}

%
% Macros for moduli stacks/spaces
%
\def\QCohstack{\mathcal{QC}\!\mathit{oh}}
\def\Cohstack{\mathcal{C}\!\mathit{oh}}
\def\Spacesstack{\mathcal{S}\!\mathit{paces}}
\def\Quotfunctor{\mathrm{Quot}}
\def\Hilbfunctor{\mathrm{Hilb}}
\def\Curvesstack{\mathcal{C}\!\mathit{urves}}
\def\Polarizedstack{\mathcal{P}\!\mathit{olarized}}
\def\Complexesstack{\mathcal{C}\!\mathit{omplexes}}
% \Pic is the operator that assigns to X its picard group, usage \Pic(X)
% \Picardstack_{X/B} denotes the Picard stack of X over B
% \Picardfunctor_{X/B} denotes the Picard functor of X over B
\def\Pic{\mathop{\mathrm{Pic}}\nolimits}
\def\Picardstack{\mathcal{P}\!\mathit{ic}}
\def\Picardfunctor{\mathrm{Pic}}
\def\Deformationcategory{\mathcal{D}\!\mathit{ef}}
\newenvironment{reference}{}{}
\newenvironment{slogan}{}{}
\newenvironment{history}{}{}
\usepackage{cleveref}
\crefname{originallemma}{original lemma}{original lemmas}
\Crefname{originallemma}{Original lemma}{Original lemmas}
\makeatletter
\newlabel{lemma-semi-local-dimension-one-conductor}{{09N2}{source}{}{sourcepointer-1}{}}
\newlabel{lemma-semi-local-dimension-one-conductor@cref}{{[originallemma][0][]09N2}{[1][0][]source}{}{}{}}
\makeatother
\crefname{originalreference}{original reference}{original references}
\Crefname{originalreference}{Original reference}{Original references}
\makeatletter
\newlabel{algebra-lemma-characterize-Dedekind}{{034X}{source}{}{sourcepointer-2}{}}
\newlabel{algebra-lemma-characterize-Dedekind@cref}{{[originalreference][0][]034X}{[1][0][]source}{}{}{}}
\makeatother
\title{Existence of an element of the base whose inverse lies in the Jacobson radicals of a finite family of one-dimensional semi-local domains}
\author{The Stacks Project Authors}
\date{stacks-project:varieties.tex:local-snapshot-20261011; GNU FDL 1.2 or later}
\begin{document}
\maketitle
\noindent\textit{Editorial note.} All delivered mathematical text is the author's own, copied verbatim from the pinned Stacks Project chapter file varieties.tex; no statement or proof step was rewritten, reordered or removed.\par\medskip
\noindent\textit{Editorial note.} The counted span cites two cross-chapter labels, each recorded in assembly.external\_references with its canonical tag as a navigation pointer.\par\medskip
\noindent\textit{Editorial note.} The proof cites no bibliography item inside the counted span.\par\medskip
\noindent\textit{Allowed standard prerequisites (editorial).}\begin{itemize}
\item Conductor of the integral closure of a one-dimensional semi-local domain. Stacks Project, Varieties, Lemma lemma-semi-local-dimension-one-conductor (tag 09N2). Allows descending from the integral closure back to B\_i after raising x to a power.
\item Characterization of Dedekind domains by discrete valuations. Stacks Project, Algebra, Lemma algebra-lemma-characterize-Dedekind (tag 034X). Supplies the discrete valuations v\_i used in the induction.
\end{itemize}\medskip
\ifdefined\hypertarget\hypertarget{sourcepointer-1}{}\fi
\noindent\textit{Source reference.} Original lemma 09N2: sources/stacks-chow/tags.txt, source lines 12200--12200. Cross-chapter standard prerequisite cited by the author inside the counted span; it is used as a pointer and its own proof is not part of this item.\par\medskip
\ifdefined\hypertarget\hypertarget{sourcepointer-2}{}\fi
\noindent\textit{Source reference.} Original reference 034X: sources/stacks-chow/tags.txt, source lines 4108--4108. Cross-chapter standard prerequisite cited by the author inside the counted span; it is used as a pointer and its own proof is not part of this item.\par\medskip
\begin{lemma}
\label{lemma-create-globally-generated}
Let $A$ be a domain with fraction field $K$. Let $B_1, \ldots, B_r \subset K$
be Noetherian $1$-dimensional semi-local domains whose fraction
fields are $K$. If $A \otimes B_i \to K$ are surjective for $i = 1, \ldots, r$,
then there exists an $x \in A$ such that $x^{-1}$ is in the Jacobson radical of
$B_i$ for $i = 1, \ldots, r$.
\end{lemma}

\begin{proof}
Let $B_i'$ be the integral closure of $B_i$ in $K$. Suppose we find a
nonzero $x \in A$ such that $x^{-1}$ is in the Jacobson radical of $B'_i$ for
$i = 1, \ldots, r$. Then by
Lemma \ref{lemma-semi-local-dimension-one-conductor},
after replacing $x$ by a power we get $x^{-1} \in B_i$.
Since $\Spec(B'_i) \to \Spec(B_i)$ is surjective we
see that $x^{-1}$ is then also in the Jacobson radical of $B_i$.
Thus we may assume that each $B_i$ is a semi-local Dedekind domain.

\medskip\noindent
If $B_i$ is not local, then remove $B_i$ from the list and
add back the finite collection of local rings $(B_i)_\mathfrak m$.
Thus we may assume that $B_i$ is a discrete valuation ring for
$i = 1, \ldots, r$.

\medskip\noindent
Let $v_i : K \to \mathbf{Z}$, $i = 1, \ldots, r$
be the corresponding discrete valuations (see
Algebra, Lemma \ref{algebra-lemma-characterize-Dedekind}).
We are looking for a nonzero $x \in A$ with $v_i(x) < 0$ for
$i = 1, \ldots, r$. We will prove this by induction on $r$.

\medskip\noindent
If $r = 1$ and the result is wrong, then $A \subset B$ and the map
$A \otimes B \to K$ is not surjective, contradiction.

\medskip\noindent
If $r > 1$, then by induction we can find a nonzero $x \in A$ such that
$v_i(x) < 0$ for $i = 1, \ldots, r - 1$. If $v_r(x) < 0$ then we are
done, so we may assume $v_r(x) \geq 0$. By the base case we can find
$y \in A$ nonzero such that $v_r(y) < 0$. After replacing $x$ by a
power we may assume that $v_i(x) < v_i(y)$ for $i = 1, \ldots, r - 1$.
Then $x + y$ is the element we are looking for.
\end{proof}

\end{document}
